% Part: normal-modal-logic
% Chapter: frame-definability
% Section: equivalence-S5

\documentclass[../../../include/open-logic-section]{subfiles}

\begin{document}

\olfileid{nml}{frd}{es5}

\olsection{સામ્ય સંબંધો અને \Log{S5}}

મોડલ તર્કશાસ્ત્ર \Log{S5} એ બધી સાર્વત્રિક ફ્રેમોમાં
માન્ય !!{formula}sના ગણ તરીકે લક્ષણિત થાય છે;
એટલે કે દરેક જગત પોતાને સહિત દરેક જગતમાંથી
સુલભ છે. આવી સ્થિતિમાં $\Box$ અનિવાર્યતાને
અને $\Diamond$ શક્યતાને અનુરૂપ છે: $!A$ \emph{દરેક}
જગતમાં સત્ય હોય તો $\Box !A$ સત્ય છે, અને $!A$
\emph{કોઈક} જગતમાં સત્ય હોય તો $\Diamond !A$
સત્ય છે. \Log{S5}ને બધી સ્વવાચક, સંમિત અને
પરંપરિત ફ્રેમોમાં, એટલે કે બધા \emph{સામ્ય સંબંધો}
પર, માન્ય !!{formula}s તરીકે પણ લક્ષણિત કરી શકાય છે.

\begin{defn}
  $W$ પરનો દ્વિસ્થાની સંબંધ $R$ \emph{સામ્ય સંબંધ}
  ત્યારે અને માત્ર ત્યારે જ છે જ્યારે તે સ્વવાચક,
  સંમિત અને પરંપરિત હોય. $W$ પરનો સંબંધ $R$
  \emph{સાર્વત્રિક} ત્યારે અને માત્ર ત્યારે જ છે જ્યારે
  બધા $u,v \in
  W$ માટે $Ruv$ હોય.
\end{defn}

\Ax{T}, \Ax{B} અને \Ax{4} અનુક્રમે સ્વવાચક,
સંમિત અને પરંપરિત ફ્રેમોને લક્ષણિત કરતાં હોવાથી,
સુલભતા સંબંધ સામ્ય સંબંધ હોય તેવી ફ્રેમો બરાબર
એ છે જેમાં આ ત્રણેય !!{formula}s માન્ય છે.
સામ્ય સંબંધોને !!{formula}sનાં અન્ય સંયોજનો વડે
પણ લક્ષણિત કરી શકાય છે, કેમ કે તેમની વ્યાખ્યામાં
વાપરેલી શરતો $R$ પરની અન્ય પરિચિત શરતોનાં
સંયોજનોને સમકક્ષ છે.

\begin{prop}\ollabel{prop:equivalences}
  નીચેની શરતો સમકક્ષ છે:
  \begin{enumerate}
  \item $R$ સામ્ય સંબંધ છે;
  \item $R$ સ્વવાચક અને યુક્લિડિયન છે;
  \item $R$ સિરિયલ, સંમિત અને યુક્લિડિયન છે;
  \item $R$ સિરિયલ, સંમિત અને પરંપરિત છે.
  \end{enumerate}
\end{prop}

\begin{proof}
  અભ્યાસપ્રશ્ન.
\end{proof}

\begin{prob}
  નીચેનું બતાવી \olref[nml][frd][es5]{prop:equivalences}
  સાબિત કરો:
  \begin{enumerate}
  \item $R$ સંમિત અને પરંપરિત હોય તો યુક્લિડિયન છે.
  \item $R$ સ્વવાચક હોય તો સિરિયલ છે.
  \item $R$ સ્વવાચક અને યુક્લિડિયન હોય તો સંમિત છે.
  \item $R$ સંમિત અને યુક્લિડિયન હોય તો પરંપરિત છે.
  \item $R$ સિરિયલ, સંમિત અને પરંપરિત હોય તો સ્વવાચક છે.
  \end{enumerate}
  આ શરતો સમકક્ષ છે તે સાબિત કરવા આ કેમ પૂરતું છે
  તે સમજાવો.
\end{prob}

\olref{prop:equivalences} એ \olref[prf][prs]{prop:S5}નો
અર્થવિજ્ઞાનલક્ષી સમકક્ષ છે: તે એવી ફ્રેમોના મોડલ
તર્કશાસ્ત્રનું સમકક્ષ લક્ષણન આપે છે જેમનો $R$
સામ્ય સંબંધ છે (જેને પરંપરાગત રીતે \Log{S5} કહે છે).

સાર્વત્રિક અને સામ્ય સંબંધો વચ્ચે શું સંબંધ છે?
દરેક સાર્વત્રિક સંબંધ સામ્ય સંબંધ છે, છતાં દરેક
સામ્ય સંબંધ સાર્વત્રિક નથી. પરંતુ બધી સાર્વત્રિક
ફ્રેમોમાં માન્ય !!{formula}s બરાબર બધી સામ્ય-સંબંધવાળી
ફ્રેમોમાં માન્ય સૂત્રો જ છે.

\begin{prop}
  $R$ સામ્ય સંબંધ માનો, અને દરેક $w \in W$ માટે
  $w$નો \emph{સામ્ય વર્ગ} એ ગણ $[w] = \{w'\in W :
  Rww'\}$ વ્યાખ્યાયિત કરો. ત્યારે:
  \begin{enumerate}
  \item $w \in [w]$;
  \item દરેક સામ્ય વર્ગ $[w]$ પર $R$ સાર્વત્રિક છે;
  \item સામ્ય વર્ગોનો સંગ્રહ $W$ને પરસ્પર વિયુક્ત
    અને મળીને સંપૂર્ણ એવા ઉપગણોમાં વિભાજિત કરે છે.
  \end{enumerate}
\end{prop}

\begin{prop}\ollabel{prop:S5=univ}
  !!^a{formula}~$!A$ એ એવી બધી ફ્રેમો
  $\mModel{F} =\tuple{W,R}$માં માન્ય છે જ્યાં $R$
  સામ્ય સંબંધ છે, ત્યારે અને માત્ર ત્યારે જ જ્યારે
  તે એવી બધી ફ્રેમો $\mModel{F} =\tuple{W,R}$માં
  માન્ય છે જ્યાં $R$ સાર્વત્રિક છે. તેથી સાર્વત્રિક
  ફ્રેમોનું તર્કશાસ્ત્ર બરાબર \Log{S5} છે.
\end{prop}

\begin{proof}
  $W$ પરનો સાર્વત્રિક સંબંધ~$R$ સામ્ય સંબંધ છે
  તે તરત ચકાસી શકાય છે. તેથી $!A$ એ $R$ સામ્ય
  સંબંધ હોય એવી બધી ફ્રેમોમાં માન્ય હોય, તો તે
  બધી સાર્વત્રિક ફ્રેમોમાં પણ માન્ય છે. બીજી
  દિશા માટે પ્રતિધનાત્મક દલીલ કરીએ: માનો કે
  !!a{formula}~$!B$ એ જગત $w$માં અસત્ય છે,
  જ્યાં નિદર્શન $\mModel{M} = \tuple{W, R, V}$
  એવી ફ્રેમ $\tuple{W,R}$ પર આધારિત છે જેમાં $R$
  એ $W$ પરનો સામ્ય સંબંધ છે. એટલે
  $\mSat/{M}{!B}[w]$. હવે નિદર્શન
  $\mModel{M}' = \tuple{W',
    R', V'}$ આ રીતે વ્યાખ્યાયિત કરો:
  \begin{enumerate}
  \item $W' = [w]$;
  \item $R'$ એ $W'$ પર સાર્વત્રિક છે;
  \item $V'(p) = V(p) \cap W'$.
  \end{enumerate}
  (અર્થાત્ \olref{fig:partition}ના ઘેરા વિસ્તારમાં
  નિદર્શન $\mModel{M}'$ના જગતોનો ગણ $W'$
  દર્શાવ્યો છે.) $W'$ પર $R$ અને $R'$ એકસરખા
  છે તે સરળતાથી જોઈ શકાય છે. પછી !!{formula}s
  પર આગમનથી બતાવી શકાય કે બધા $w' \in W'$ માટે:
  દરેક $!A$ માટે $\mSat{M'}{!A}[w']$ ત્યારે અને
  માત્ર ત્યારે જ જ્યારે $\mSat{M}{!A}[w']$ હોય
  ($W'
  \subseteq W$ હોવાથી આ વિધાન અર્થપૂર્ણ છે).
  ખાસ કરીને $\mSat/{M'}{!B}[w]$, એટલે $!B$
  સાર્વત્રિક ફ્રેમ પર આધારિત નિદર્શનમાં અસત્ય છે.
\end{proof}

\begin{figure}[t]
  \centering
  \begin{tikzpicture}[node distance=2cm, auto, thick]
    \clip [rounded corners] (0,0) -- (6,0) -- (6,4) -- (0,4) --  cycle;
    \begin{scope}
      \clip (0,0) -- (2,0)  .. controls (1.5,1.5) and (4.5,1.5)
      .. (4,4) -- (0,4) -- (0,0) -- cycle;
      \filldraw[fill=gray!40] (0,4) -- (0,2) .. controls (1.5,1.5)
      and (4.5,1.5) .. (4.5,0) -- (6,0) -- (6,4) -- (0,4) --  cycle;
    \end{scope}
    \draw [name path=line1] (2,0) .. controls (1.5,1.5) and (4.5,1.5) .. (4,4);
    \draw [name path=line2] (0,2)  .. controls (1.5,1.5) and (4.5,1.5) .. (4.5,0);
    \path [name intersections={of=line1 and line2}];
    \draw [rounded corners, use as bounding box, very thick] (0,0)
    -- (6,0) -- (6,4) -- (0,4) --  cycle; 
    \node at (2,3) {$[w]$} ;
    \node at (1,1) {$[u]$} ;
    \node at (3,0.75) {$[v]$} ;
    \node at (4.75,2) {$[z]$} ;
  \end{tikzpicture}
  \caption{સામ્ય વર્ગોમાં $W$નું વિભાજન.}
  \ollabel{fig:partition}
\end{figure}

\end{document}
